% genere par scripts/rechnung_retraite.py (quellen_tex) - ne pas modifier a la main
\expandafter\newcommand\csname QHfn@abgeltungsteuer_satz\endcsname{§32d Abs. 1 EStG. \href{https://www.gesetze-im-internet.de/estg/__32d.html}{gesetze-im-internet.de}, consulté le 29.09.2026 (source primaire).}
\expandafter\newcommand\csname QHfn@basiszins_2026\endcsname{BMF-Schreiben vom 13.01.2026, Basiszins zur Berechnung der Vorabpauschale zum 2. Januar 2026, IV C 1 - S 1980/00230/012/001. \href{https://www.bundesfinanzministerium.de/Content/DE/Downloads/BMF_Schreiben/Steuerarten/Investmentsteuer/2026-01-13-basiszins-berechnung-vorabpauschale.html}{bundesfinanzministerium.de}, consulté le 29.09.2026 (source secondaire).}
\expandafter\newcommand\csname QHfn@durchschnittsentgelt_2026\endcsname{§3 Abs. 2 Sozialversicherungsrechengroessen-Verordnung 2026 (SVBezGrV 2026). \href{https://www.gesetze-im-internet.de/svbezgrv_2026/__3.html}{gesetze-im-internet.de}, consulté le 29.09.2026 (source primaire).}
\expandafter\newcommand\csname QHfn@einstiegsgehalt_brutto\endcsname{StepStone Gehaltsreport 2026, Medianwert Ingenieurinnen/Ingenieure mit unter einem Jahr Berufserfahrung (Datenbasis Jan. 2022 - Nov. 2025). \href{https://www.stepstone.de/e-recruiting/hr-wissen/gehalt/gehaltsreport-deutschlands-gehaelter-im-fokus/}{stepstone.de}, consulté le 29.09.2026 (source secondaire).}
\expandafter\newcommand\csname QHfn@est_tarif_2026\endcsname{§32a Abs. 1 EStG, Veranlagungszeitraum 2026. \href{https://www.gesetze-im-internet.de/estg/__32a.html}{gesetze-im-internet.de}, consulté le 29.09.2026 (source primaire).}
\expandafter\newcommand\csname QHfn@etf_welt_rendite_div\endcsname{Repli explicite (Annahme), aucune valeur exploitable trouvee sur la fiche justETF du plus grand ETF MSCI World. \href{https://www.justetf.com/en/etf-profile.html?isin=IE00B4L5Y983}{justetf.com}, consulté le 30.09.2026 (source secondaire).}
\expandafter\newcommand\csname QHfn@gehaltssteigerung_real\endcsname{Abgeleitet aus StepStone Gehaltsreport 2026 (Querschnitt nach Berufserfahrung: Einstieg 59.250 EUR vs. 25 Jahre Berufserfahrung 93.250 EUR, impliziert ca. 1,8 \%/Jahr Wachstumsrate). \href{https://www.stepstone.de/e-recruiting/hr-wissen/gehalt/gehaltsreport-deutschlands-gehaelter-im-fokus/}{stepstone.de}, consulté le 29.09.2026 (source secondaire).}
\expandafter\newcommand\csname QHfn@inflation_2022_de\endcsname{Statistisches Bundesamt (Destatis), Pressemitteilung Nr. 022 vom 17.01.2023, 'Inflationsrate im Jahr 2022 bei +7,9 \%'. \href{https://www.destatis.de/DE/Presse/Pressemitteilungen/2023/01/PD23_022_611.html}{destatis.de}, consulté le 29.09.2026 (source primaire).}
\expandafter\newcommand\csname QHfn@inflation_destatis_2019_2025\endcsname{Statistisches Bundesamt (Destatis), Pressemitteilungen annuelles 'Inflationsrate im Jahr JJJJ': PD20\_019\_611 (2019), PD21\_025\_611 (2020), PD22\_025\_611 (2021), PD25\_020\_611 (citant retrospectivement 2022 a +6,9\% et 2023 a +5,9\%), PD24\_020\_611 (2023), PD25\_020\_611 (2024), PD26\_019\_611 KORREKTUR (2025). \href{https://www.destatis.de/DE/Presse/Pressemitteilungen/2025/01/PD25_020_611.html}{destatis.de}, consulté le 30.09.2026 (source primaire).}
\expandafter\newcommand\csname QHfn@inflation_ziel\endcsname{EZB, geldpolitische Strategie, Preisstabilitaetsziel von 2 \% (HVPI). \href{https://www.ecb.europa.eu/mopo/strategy/pricestab/html/index.de.html}{ecb.europa.eu}, consulté le 29.09.2026 (source primaire).}
\expandafter\newcommand\csname QHfn@kv_bbg_monat_2026\endcsname{GKV-Spitzenverband-Faktenblatt 'Rechengroessen und Grenzwerte im Versicherungs- und Beitragsrecht fuer das Jahr 2026', Beitragsbemessungsgrenzen Kranken- und Pflegeversicherung, S. 1. \href{https://www.gkv-spitzenverband.de/media/dokumente/presse/zahlen_und_grafiken/20260101_Faktenblatt_Rechengroessen_Beitragsrecht.pdf}{gkv-spitzenverband.de}, consulté le 29.09.2026 (source primaire).}
\expandafter\newcommand\csname QHfn@kv_mindestbemessung_monat_2026\endcsname{GKV-Spitzenverband-Faktenblatt 'Rechengroessen und Grenzwerte im Versicherungs- und Beitragsrecht fuer das Jahr 2026', Freiwillige Versicherung in der Krankenversicherung, Mindestbemessungsgrundlage, S. 1. \href{https://www.gkv-spitzenverband.de/media/dokumente/presse/zahlen_und_grafiken/20260101_Faktenblatt_Rechengroessen_Beitragsrecht.pdf}{gkv-spitzenverband.de}, consulté le 29.09.2026 (source primaire).}
\expandafter\newcommand\csname QHfn@kv_satz_ermaessigt\endcsname{§243 SGB V; bestaetigt im GKV-Spitzenverband-Faktenblatt 'Rechengroessen und Grenzwerte im Versicherungs- und Beitragsrecht fuer das Jahr 2026' (1. Januar 2026), Beitragssaetze, S. 1. \href{https://www.gkv-spitzenverband.de/media/dokumente/presse/zahlen_und_grafiken/20260101_Faktenblatt_Rechengroessen_Beitragsrecht.pdf}{gkv-spitzenverband.de}, consulté le 29.09.2026 (source primaire).}
\expandafter\newcommand\csname QHfn@kv_schwellen_dynamisierung\endcsname{§223 Abs. 4 SGB V: 'Die Beitragsbemessungsgrenze im Jahr 2027 entspricht der um 3 600 Euro erhoehten Jahresarbeitsentgeltgrenze nach § 6 Absatz 7 [SGB V] ... Die Bundesregierung setzt die Beitragsbemessungsgrenze in der Rechtsverordnung nach § 160 des Sechsten Buches gesondert fest.' §6 Abs. 6 Satz 2 SGB V (Jahresarbeitsentgeltgrenze, via Abs. 7 referenziert): 'Sie aendert sich zum 1. Januar eines jeden Jahres in dem Verhaeltnis, in dem die Bruttoloehne und -gehaelter je Arbeitnehmer ... im vergangenen Kalenderjahr zu den entsprechenden Bruttoloehnen und -gehaeltern im vorvergangenen Kalenderjahr stehen', d.h. eine jaehrliche Lohn-Indexierung, keine feste nominale Grenze. Die Mindestbemessungsgrundlage fuer freiwillig Versicherte ist ein Bruchteil der Bezugsgroesse (§18 SGB IV): 'das Durchschnittsentgelt der gesetzlichen Rentenversicherung im vorvergangenen Kalenderjahr, aufgerundet auf den naechsthoeheren, durch 420 teilbaren Betrag', ebenfalls jaehrlich per Sozialversicherungs-Rechengroessenverordnung festgesetzt.. \href{https://www.gesetze-im-internet.de/sgb_5/__223.html}{gesetze-im-internet.de}, consulté le 30.09.2026 (source primaire).}
\expandafter\newcommand\csname QHfn@kv_teilfreistellung_etf\endcsname{GKV-Spitzenverband, 'Katalog von Einnahmen und deren beitragsrechtliche Bewertung nach § 240 SGB V', Stand 26.05.2026, Zeile 'Investmenterträge': '§ 16 InvStG | ja, unter Beruecksichtigung der §§ 20 und 56 Abs. 6 InvStG'. Die separate Zeile 'Kapitalvermoegen, Einkuenfte aus -' (individuelle Aktiendividenden, kein Fondsvehikel) traegt keinen solchen Verweis auf §20 InvStG., 15 von 30, Zeile 'Investmenterträge'. \href{https://www.gkv-spitzenverband.de/media/dokumente/krankenversicherung_1/grundprinzipien_1/finanzierung/beitragsbemessung/2026-05-26_Katalog_Einnahmen_beitragsrechtliche_Bewertung_240_SGB_V_BF.pdf}{gkv-spitzenverband.de}, consulté le 30.09.2026 (source primaire).}
\expandafter\newcommand\csname QHfn@kv_zusatzbeitrag_2026\endcsname{GKV-Spitzenverband-Faktenblatt 'Rechengroessen und Grenzwerte im Versicherungs- und Beitragsrecht fuer das Jahr 2026', Beitragssaetze, Beitragssaetze, S. 1. \href{https://www.gkv-spitzenverband.de/media/dokumente/presse/zahlen_und_grafiken/20260101_Faktenblatt_Rechengroessen_Beitragsrecht.pdf}{gkv-spitzenverband.de}, consulté le 29.09.2026 (source primaire).}
\expandafter\newcommand\csname QHfn@pv_satz_kinderlos_2026\endcsname{GKV-Spitzenverband-Faktenblatt 'Rechengroessen und Grenzwerte im Versicherungs- und Beitragsrecht fuer das Jahr 2026', Beitragssaetze (3,6 \% + 0,6 \% Zuschlag fuer Kinderlose), Beitragssaetze, S. 1. \href{https://www.gkv-spitzenverband.de/media/dokumente/presse/zahlen_und_grafiken/20260101_Faktenblatt_Rechengroessen_Beitragsrecht.pdf}{gkv-spitzenverband.de}, consulté le 29.09.2026 (source primaire).}
\expandafter\newcommand\csname QHfn@quellensteuer\endcsname{Bundeszentralamt fuer Steuern (BZSt), 'Anrechenbarkeit der Quellensteuer auf Dividenden und Zinsen von Staaten, mit denen Deutschland ein Doppelbesteuerungsabkommen abgeschlossen hat', Stand 1. Januar 2025, Redaktionsschluss Juni 2025, laenderspezifische Zeilen, S. 1-25 (Belgien S.5, Daenemark S.7, Finnland S.9, Frankreich S.9, Irland S.12, Italien S.13, Kanada S.14, Niederlande S.21, Norwegen S.22, Spanien S.24, Schweiz S.23, Vereinigte Staaten S.24, Vereinigtes Koenigreich S.24). \href{https://www.bzst.de/SharedDocs/Downloads/DE/EU_OECD/anrechenbare_ausl_quellensteuer_2025.pdf}{bzst.de}, consulté le 29.09.2026 (source primaire).}
\expandafter\newcommand\csname QHfn@regelaltersgrenze\endcsname{§35 SGB VI, Regelaltersgrenze fuer Jahrgaenge ab 1964. \href{https://www.deutsche-rentenversicherung.de/DRV/DE/Experten/Arbeitgeber-und-Steuerberater/summa-summarum/Lexikon/R/regelaltersgrenze}{deutsche-rentenversicherung.de}, consulté le 29.09.2026 (source primaire).}
\expandafter\newcommand\csname QHfn@rentenwert_2026\endcsname{Deutsche Rentenversicherung, Pressemitteilung 'Bundeskabinett beschliesst Rentenanpassung 2026' (29.04.2026), Anstieg von 40,79 auf 42,52 EUR. \href{https://www.deutsche-rentenversicherung.de/DRV/DE/Ueber-uns-und-Presse/Presse/Meldungen/2026/260429-rentenanpassung-2026-bundeskabinett.html}{deutsche-rentenversicherung.de}, consulté le 29.09.2026 (source primaire).}
\expandafter\newcommand\csname QHfn@shiller_url\endcsname{Robert J. Shiller, Yale University, Online-Datentabelle S\&P 500 (Kurse, Dividenden, Ertraege, CPI) seit 1871. \href{http://www.econ.yale.edu/~shiller/data/ie_data.xls}{econ.yale.edu}, consulté le 29.09.2026 (source primaire).}
\expandafter\newcommand\csname QHfn@soli_freigrenze_2026\endcsname{§3 Abs. 3 SolzG 1995 n.F. (Steuerfortentwicklungsgesetz 2024), erstmals VZ 2026. \href{https://www.buzer.de/gesetz/282/al234205-0.htm}{buzer.de}, consulté le 29.09.2026 (source primaire).}
\expandafter\newcommand\csname QHfn@soli_satz\endcsname{§4 SolzG 1995. \href{https://www.gesetze-im-internet.de/solzg_1995/__4.html}{gesetze-im-internet.de}, consulté le 29.09.2026 (source primaire).}
\expandafter\newcommand\csname QHfn@sparerpauschbetrag\endcsname{§20 Abs. 9 EStG, seit 2023. \href{https://www.gesetze-im-internet.de/estg/__20.html}{gesetze-im-internet.de}, consulté le 29.09.2026 (source primaire).}
\expandafter\newcommand\csname QHfn@sv_arbeitnehmer\endcsname{GKV-Spitzenverband-Faktenblatt 'Rechengroessen und Grenzwerte im Versicherungs- und Beitragsrecht fuer das Jahr 2026', S. 1. \href{https://www.gkv-spitzenverband.de/media/dokumente/presse/zahlen_und_grafiken/20260101_Faktenblatt_Rechengroessen_Beitragsrecht.pdf}{gkv-spitzenverband.de}, consulté le 29.09.2026 (source primaire).}
\expandafter\newcommand\csname QHfn@teilfreistellung_aktienfonds\endcsname{§20 InvStG (Aktienfonds, natürliche Person im Privatvermögen). \href{https://www.gesetze-im-internet.de/invstg_2018/__20.html}{gesetze-im-internet.de}, consulté le 29.09.2026 (source primaire).}
